% ============================================================================
% 第 2 章 解决方法概要
% ============================================================================
\section{解决方法概要}
\label{sec:methodology}

\subsection{三层协同调度架构}

为破解任务指派与时空路径规划之间的强耦合矛盾，同时兼顾大规模仓储环境下的计算效率与物理安全性，本系统采用\textbf{“上层任务指派—中层时空避碰—底层独立验解”}的三层协同架构（如图~\ref{fig:system-arch} 所示）。整个在线调度循环严格遵循因果信息边界：在每个离散时间步，调度器仅接收当前已发布且未完成的可见订单、车辆实时位姿及已承诺的未来时空预约表，输出原子可执行的调度计划并交由统一仿真内核步进推进。

\begin{figure}[htbp]
  \centering
  \begin{tikzpicture}[
    node distance=0.65cm and 0.9cm,
    box/.style={
      rectangle, rounded corners=3pt, draw=PrimaryNavy, thick,
      align=center, font=\small, inner sep=6pt, minimum height=1.15cm
    },
    layer/.style={
      rectangle, rounded corners=4pt, draw=BorderGray, dashed, thick,
      inner sep=8pt
    },
    arr/.style={-{Stealth[length=2.4mm]}, thick, color=PrimaryNavy},
    darr/.style={-{Stealth[length=2.2mm]}, thick, dashed, color=AccentTeal}
  ]
    % 顶部输入
    \node[box, fill=SoftBlueBg, minimum width=3.8cm] (env) {
      \sffamily\bfseries 仓库环境与动态订单流 \\
      \footnotesize 布局几何 / 层序状态 / 已发布订单
    };

    % 三层核心
    \node[box, fill=white, draw=AccentBlue, below=0.85cm of env, minimum width=4.3cm] (dispatch) {
      \sffamily\bfseries\color{PrimaryNavy} 1. 任务指派与决策层 \\
      \footnotesize 贪心 / 匹配 / 影子试派（BCRD）\\
      \footnotesize 事件门控 + 软分区联合 GA（\code{gate}）
    };

    \node[box, fill=white, draw=AccentTeal, below=0.85cm of dispatch, minimum width=4.3cm] (routing) {
      \sffamily\bfseries\color{AccentTeal} 2. 时空无碰撞路径规划层 \\
      \footnotesize 增量时空预约表 + SIPP 安全区间规划 \\
      \footnotesize 物理预演评分 / 拥堵关联 LNS 联合修复
    };

    \node[box, fill=SoftTealBg, draw=AccentTeal, below=0.85cm of routing, minimum width=4.3cm] (commit) {
      \sffamily\bfseries\color{PrimaryNavy} 3. 唯一仿真状态机与原子提交 \\
      \footnotesize 全未来冲突审计 $\rightarrow$ 步进推进 $\rightarrow$ 返库收尾
    };

    % 右侧独立验解与评估
    \node[box, fill=SoftAmberBg, draw=WarmAmber, right=1.35cm of routing, minimum width=3.5cm] (verify) {
      \sffamily\bfseries\color{WarmAmber} 独立物理验解器 \\
      \footnotesize 顶点/换边无碰撞校验 \\
      \footnotesize 装卸时长/层序/返库校验
    };

    \node[box, fill=SoftBlueBg, draw=PrimaryNavy, below=0.85cm of verify, minimum width=3.5cm] (metrics) {
      \sffamily\bfseries 综合指标与可视化 \\
      \footnotesize 指标 $T, R, D, F$ 统计 \\
      \footnotesize Web 回放 / 敏感性分析
    };

    % 连线
    \draw[arr] (env) -- node[right, font=\footnotesize, color=MutedText]{因果可见观测} (dispatch);
    \draw[arr] (dispatch) -- node[right, font=\footnotesize, color=MutedText]{候选任务分配与序列} (routing);
    \draw[arr] (routing) -- node[right, font=\footnotesize, color=MutedText]{无碰撞时空轨迹} (commit);

    % 反馈闭环（左侧）
    \draw[arr] (commit.west) -- ++(-1.15, 0) |- node[pos=0.25, left, font=\footnotesize, align=right, color=MutedText]{状态更新\\预约表滚动} (env.west);

    % 规划层与指派层双向预演反馈
    \draw[darr] (routing.east) -- ++(0.45, 0) |- node[pos=0.25, right, font=\footnotesize, color=AccentTeal]{预演代价} (dispatch.east);

    % 输出到独立验解
    \draw[arr] (commit.east) -| ($(verify.west)+(-0.35,0)$) |- (verify.west);
    \draw[arr] (verify) -- node[right, font=\footnotesize, color=MutedText]{100\% 可行确认} (metrics);
  \end{tikzpicture}
  \caption{多 AGV 仓库智能调度系统三层协同与独立验解架构图}
  \label{fig:system-arch}
\end{figure}

三层架构的核心职责划分如下：

\begin{itemize}
  \item \textbf{任务指派与决策层}：决定“由哪辆车以何种顺序响应哪些就绪任务”。在低负载、少量并发场景下采用低开销精确或贪心策略直接派单；在高负载或强竞争场景下，通过滚动影子试派（主线 A）或事件门控触发的软分区联合遗传搜索（主线 B）在多车多单组合空间内寻找全局协调解。
  \item \textbf{时空无碰撞路径规划层}：维护所有已指派车辆的离散时空预约表（记录顶点占用与有向边占用），并基于\textbf{安全区间路径规划（SIPP）}~\cite{phillips2011sipp} 将连续空闲时间段压缩为最大安全区间，从而在显著缩减状态搜索空间的同时，直接处理货架连续装卸占位、交付站即卸即走及窄巷道避让。此外，规划层一方面为上层候选方案提供真实的“预约表物理预演代价”，另一方面在主线 A 中支持对局部拥堵车队执行大邻域搜索（LNS）联合后缀重排~\cite{li2021mapflns}。
  \item \textbf{唯一仿真状态机与独立物理验解层}：所有规划结果在提交前均须通过全未来路径冲突与服务连续性审计，采用原子提交方式更新系统状态；仿真结束后，由与调度算法完全解耦的独立验解模块逐帧核查顶点冲突、相向换边冲突、装卸持续时长、同面货架层序解锁以及最终返库状态，确保实验数据绝对可信。
\end{itemize}

\subsection{双主线研究布局与最终采用方案}

考虑到课程项目既需要具备清晰通用、可扩展交互演示能力的工程平台，又需要在贴近工业界真实到达分布的大规模数据上验证核心算法的泛化性能，我们将整体研发组织为两条互补的研究主线，并明确了各自的最终定型版本与代码入口（见表~\ref{tab:dual-track}）：

\begin{table}[htbp]
  \centering
  \caption{本项目双研究主线定位、最终采用方案与核心机制总览}
  \label{tab:dual-track}
  \small
  \renewcommand{\arraystretch}{1.28}
  \begin{tabularx}{\textwidth}{>{\raggedright\arraybackslash}p{2.35cm} >{\raggedright\arraybackslash}p{3.0cm} >{\raggedright\arraybackslash}X >{\raggedright\arraybackslash}p{3.35cm}}
    \toprule
    \sffamily\bfseries 研究主线 & \sffamily\bfseries 最终采用方案 & \sffamily\bfseries 核心算法与关键机制 & \sffamily\bfseries 源码模块 \\
    \midrule
    \textbf{主线 B}：真实工业订单流（核心实验主线） &
    WEPA 推荐版 \code{gate}（算力受限时可选 \code{gate\_preview3}） &
    面向 $62\times 47$ 地图、6 车、5 层、900\,s 真实订单窗口：采用事件门控滚动派单（低负载精确分支/免触发 GA）、软分区联合 GA 寻优，并结合预约表 SIPP 物理预演择优提交。 &
    \footnotesize\code{research/}\newline\code{recommended-20261024/} \\
    \midrule
    \textbf{主线 A}：自研通用调度与可视化平台 &
    \code{src/agv/} v0.6.0（V9.1 + SIPP/BCRD 引擎） &
    支持 Legacy、V4 分层单面与 V9/V9.1 混合背靠背货架：V4 分层动态采用 \code{bcrd\_budget}（影子试派）；Legacy 静态采用 \code{sipp\_batch}；V9.1 静态采用 \code{sipp\_stream\_bcrd}（\code{max\_batches=2}）；内置本地 Web 回放。 &
    \footnotesize\code{src/agv/}（CLI 入口 \code{agv} 与 \code{agv serve}） \\
    \bottomrule
  \end{tabularx}
\end{table}

两条主线在底层共享同一套离散时空无碰撞建模思想与 SIPP 路径核心，但在实验侧重上形成良好互补：主线 A 侧重\textbf{多工况抽象与机理验证}，覆盖从规则栅格到单面多层、背靠背混合货架及流式关闭的渐进建模，并支持教师与评测人员通过浏览器直观观察 AGV 取货、避碰与返库过程；主线 B 则聚焦\textbf{真实工业数据驱动的深度算法优化}，依托公开的 WEPAStacks 真实入出库请求流（411,830 条原始记录抽样窗口）开展多方案配对比较、开发/验证集分离泛化测试与多随机种子敏感性分析，最终收敛出大作业核心推荐方案 \code{gate}。

\subsection{统一评价口径与工程复现保障}

为避免单一指标失衡，两条主线均采用兼顾周期效率、单单响应与能耗里程的加权综合目标函数 $F$（数值越低越优）：
\begin{itemize}
  \item \textbf{主线 B（独立物理尺度线性评分）}：去除对特定基线算法分母的依赖，按 900\,s 窗口时长与 6 车总步长尺度固定归一化，定义为
  \begin{equation}
    F_{\text{WEPA}} = 0.5\frac{T}{900} + 0.3\frac{R}{900} + 0.2\frac{D}{5400}
    \label{eq:obj-wepa-overview}
  \end{equation}
  其中 $T$ 为最后一笔订单的完成时刻（不含最终返库阶段），$R$ 为全部订单从发布到完成的平均耗时，$D$ 为含最终返库在内的全车队实际移动总步数。
  \item \textbf{主线 A（同场景参照归一化评分）}：在 V4 与 V9.1 等多尺度合成场景下，以同场景固定时空 A* 贪心基线指标 $(T_{\mathrm{ref}}, R_{\mathrm{ref}}, D_{\mathrm{ref}})$ 为参照基准 $1.0$，定义为 $F_{\text{Core}} = 0.5(T/T_{\mathrm{ref}}) + 0.3(R/R_{\mathrm{ref}}) + 0.2(D/D_{\mathrm{ref}})$；在 Legacy 场景下则直接对比平均流程时间 $R$、完工时间 $T$ 与总步数 $D$。
\end{itemize}

在工程实现方面，两条主线的核心调度、仿真与验解模块均仅依赖 \textbf{Python 3.10+ 标准库}，无需安装商业优化求解器或重型深度学习框架即可在跨平台环境下直接运行。同时，仓库内固化了完整的单元与集成测试套件、冻结证据校验脚本以及可重复的命令行复现入口，为后续章节的定量实验提供了严格的可审计基础。
